% ============================================================
% FÓRMULAS — CARGA ELÉCTRICA Y ELECTRIZACIÓN
% ============================================================

\section{Fórmulas: carga eléctrica y electrización}


% ============================================================
% CARGA ELEMENTAL
% ============================================================

\subsection{Carga elemental}

La magnitud de la carga elemental es

\begin{equation}
    \boxed{
    e=\SI{1.602176634e-19}{\coulomb}
    }.
\end{equation}

Para las principales partículas del átomo:

\begin{equation}
    \boxed{
    q_p=+e
    }
\end{equation}

\begin{equation}
    \boxed{
    q_e=-e
    }
\end{equation}

\begin{equation}
    \boxed{
    q_n=0
    }.
\end{equation}

donde

\begin{align*}
    q_p &= \text{carga del protón},\\
    q_e &= \text{carga del electrón},\\
    q_n &= \text{carga neta del neutrón}.
\end{align*}


% ============================================================
% CARGA TOTAL
% ============================================================

\subsection{Carga total de un sistema}

La carga eléctrica es una magnitud escalar y se suma algebraicamente:

\begin{equation}
    \boxed{
    Q=q_1+q_2+\cdots+q_n
    }.
\end{equation}

De forma compacta,

\begin{equation}
    \boxed{
    Q=\sum_i q_i
    }.
\end{equation}

donde

\begin{align*}
    Q   &= \text{carga total del sistema},\\
    q_i &= \text{cada una de las cargas que forman el sistema}.
\end{align*}


% ============================================================
% CUERPO NEUTRO
% ============================================================

\subsection{Condición de neutralidad}

Un cuerpo eléctricamente neutro cumple

\begin{equation}
    \boxed{
    Q=0
    }.
\end{equation}

Equivalentemente,

\begin{equation}
    \boxed{
    Q_{+}+Q_{-}=0
    }.
\end{equation}

Esto significa que la carga positiva total y la carga negativa total se
compensan.


% ============================================================
% ELECTRONES TRANSFERIDOS
% ============================================================

\subsection{Carga producida por transferencia de electrones}

Si un cuerpo inicialmente neutro \textbf{pierde} \(N\) electrones,

\begin{equation}
    \boxed{
    Q=+Ne
    }.
\end{equation}

Por tanto,

\begin{equation}
    \boxed{
    \text{pierde electrones}
    \Longrightarrow
    Q>0
    }.
\end{equation}

Si un cuerpo inicialmente neutro \textbf{gana} \(N\) electrones,

\begin{equation}
    \boxed{
    Q=-Ne
    }.
\end{equation}

Por tanto,

\begin{equation}
    \boxed{
    \text{gana electrones}
    \Longrightarrow
    Q<0
    }.
\end{equation}


% ============================================================
% NÚMERO DE ELECTRONES
% ============================================================

\subsection{Número de electrones transferidos}

Si se conoce la magnitud de la carga adquirida por un cuerpo,

\begin{equation}
    \boxed{
    N=\frac{\abs{Q}}{e}
    }.
\end{equation}

donde

\begin{align*}
    N &= \text{número de electrones transferidos},\\
    Q &= \text{carga neta adquirida},\\
    e &= \text{carga elemental}.
\end{align*}

El número \(N\) debe ser entero.

El signo de \(Q\) indica qué ocurrió:

\begin{equation}
    \boxed{
    Q>0
    \Longrightarrow
    \text{el cuerpo perdió electrones}
    }
\end{equation}

\begin{equation}
    \boxed{
    Q<0
    \Longrightarrow
    \text{el cuerpo ganó electrones}
    }.
\end{equation}


% ============================================================
% CONSERVACIÓN DE LA CARGA
% ============================================================

\subsection{Conservación de la carga eléctrica}

Para un sistema aislado,

\begin{equation}
    \boxed{
    Q_{\text{total,inicial}}
    =
    Q_{\text{total,final}}
    }.
\end{equation}

De manera equivalente,

\begin{equation}
    \boxed{
    \sum Q_{\text{inicial}}
    =
    \sum Q_{\text{final}}
    }.
\end{equation}

Si dos cuerpos inicialmente neutros intercambian carga y el sistema permanece
aislado,

\begin{equation}
    Q_1+Q_2=0,
\end{equation}

por lo que

\begin{equation}
    \boxed{
    Q_1=-Q_2
    }.
\end{equation}


% ============================================================
% CUANTIZACIÓN
% ============================================================

\subsection{Cuantización de la carga eléctrica}

La carga neta aparece en múltiplos enteros de la carga elemental:

\begin{equation}
    \boxed{
    Q=ne
    }.
\end{equation}

con

\begin{equation}
    \boxed{
    n=0,\pm1,\pm2,\pm3,\ldots
    }.
\end{equation}

donde

\begin{align*}
    Q &= \text{carga eléctrica neta},\\
    n &= \text{número entero de cargas elementales},\\
    e &= \SI{1.602176634e-19}{\coulomb}.
\end{align*}


% ============================================================
% EQUIVALENCIA DEL COULOMB
% ============================================================

\subsection{Cantidad de cargas elementales en un coulomb}

El número de cargas elementales correspondiente a una carga de magnitud
\(\SI{1}{\coulomb}\) es

\begin{equation}
    N
    =
    \frac{\SI{1}{\coulomb}}{e}.
\end{equation}

Por tanto,

\begin{equation}
    \boxed{
    N\approx6.24\times10^{18}
    }.
\end{equation}

Así,

\begin{equation}
    \boxed{
    \SI{1}{\coulomb}
    \approx
    6.24\times10^{18}
    \text{ cargas elementales}
    }.
\end{equation}


% ============================================================
% POLARIZACIÓN
% ============================================================

\subsection{Polarización de un cuerpo neutro}

Un cuerpo puede estar polarizado y continuar siendo eléctricamente neutro:

\begin{equation}
    \boxed{
    Q_{\text{neto}}=0
    }.
\end{equation}

Por tanto,

\[
\boxed{
\text{polarización}
\not\Rightarrow
\text{carga neta distinta de cero}
}
\]

La polarización corresponde a una redistribución espacial de las cargas.


% ============================================================
% ELECTRIZACIÓN POR FROTAMIENTO
% ============================================================

\subsection{Electrización por frotamiento}

Si dos cuerpos inicialmente neutros intercambian electrones y forman un
sistema aislado,

\begin{equation}
    Q_{\text{total,inicial}}=0
\end{equation}

y después de la transferencia,

\begin{equation}
    Q_1+Q_2=0.
\end{equation}

Por tanto,

\begin{equation}
    \boxed{
    Q_1=-Q_2
    }.
\end{equation}

Los cuerpos adquieren cargas de igual magnitud y signos opuestos, siempre que
se considere el sistema formado por ambos cuerpos y no exista intercambio de
carga con el exterior.


% ============================================================
% ELECTRIZACIÓN POR INDUCCIÓN
% ============================================================

\subsection{Electrización por inducción}

Durante la primera etapa de la inducción, al acercar el inductor sin permitir
todavía transferencia de carga,

\begin{equation}
    \boxed{
    Q_{\text{conductor}}=0
    }.
\end{equation}

El conductor está polarizado, pero todavía es neutro.

Después de completar correctamente el proceso de inducción:

\[
\boxed{
\text{inductor negativo}
\Longrightarrow
\text{cuerpo inducido positivo}
}
\]

y

\[
\boxed{
\text{inductor positivo}
\Longrightarrow
\text{cuerpo inducido negativo}
}.
\]

